% Part: first-order-logic 
% Chapter: axiomatic-deduction 
% Section: deduction-theorem-quantifiers

\documentclass[../../../include/open-logic-section]{subfiles}

\begin{document}

\olfileid{fol}{axd}{ddq}

\olsection{De deductiestelling met kwantoren}

\begin{thm}[Deductiestelling]
\ollabel{thm:deduction-thm-q} Als $\Gamma \cup \{!A\} \Proves !B$,
dan $\Gamma \Proves !A \lif !B$.
\end{thm}

\begin{proof}
We gebruiken opnieuw inductie naar de lengte van de !!{derivation}
van $!B$ uit $\Gamma \cup \{!A\}$.

Het bewijs van de inductiebasis is identiek aan dat in het bewijs
van~\olref[ded]{thm:deduction-thm}.

Veronderstel voor de inductiestap opnieuw dat de !!{derivation} van
$!B$ uit $\Gamma \cup \{!A\}$ eindigt met een stap~$!B$ die door een
afleidingsregel wordt gerechtvaardigd. Als die afleidingsregel modus
ponens is, gaan we te werk als in het bewijs van
\olref[ded]{thm:deduction-thm}. Als die afleidingsregel \QR{} is, weten
we dat $!B \ident !C \lif \lforall[x][!D(x)]$ en dat eerder in de
!!{derivation} !!a{formula} van de vorm $!C \lif !D(a)$ voorkomt,
waarbij $a$ niet voorkomt in~$!C$, $!A$ of $\Gamma$. Dus geldt
\begin{align*}
  \Gamma \cup \{!A\} & \Proves !C \lif !D(a),\\
  \intertext{en is de inductiehypothese van toepassing, dat wil zeggen}
    \Gamma & \Proves !A \lif (!C \lif !D(a)).\\
  \intertext{Volgens}
  & \Proves (!A \lif (!C \lif !D(a))) \lif ((!A \land !C) \lif !D(a))\\
  \intertext{en modus ponens krijgen we}
  \Gamma & \Proves (!A \land !C) \lif !D(a).\\
  \intertext{Omdat de eigenvariabelevoorwaarde nog steeds geldt, kunnen we aan deze !!{derivation} een door \QR{} gerechtvaardigde stap toevoegen en krijgen we}
    \Gamma & \Proves (!A \land !C) \lif \lforall[x][!D(x)].\\
    \intertext{We hebben ook}
    & \Proves ((!A \land !C) \lif \lforall[x][!D(x)]) \lif (!A \lif (!C \lif \lforall[x][!D(x)]),\\
    \intertext{zodat volgens modus ponens}
    \Gamma & \Proves !A \lif (!C \lif \lforall[x][!D(x)]),
\end{align*}
dat wil zeggen, $\Gamma \Proves !B$.

We laten het geval waarin $!B$ door de regel \QR{} wordt
gerechtvaardigd maar de vorm $\lexists[x][!D(x)] \lif !C$ heeft, als
oefening.
\end{proof}

\begin{prob}
  Voltooi het bewijs van \olref[fol][axd][ddq]{thm:deduction-thm-q}.
\end{prob}

\end{document}
